\documentclass{praxisplaus}

\journalvolume{32}
\journalissue{4}
\journalyear{2026}
\articledate{12.10.2026}
\articlepages{201--209}

\title{Zur praktischen Verwendbarkeit hinreichender Gründe}

\subtitle{Bemerkungen zu einer häufig unterschätzten Entscheidungshilfe}

\author{Dr. Maximilian Methodius}

\shorttitlecommand{Hinreichende Gründe}

\articlekeywords{
  Entscheidung; Begründung; Plausibilität; Alltag
}

\abstracttext{
Der Beitrag untersucht die Frage, wann ein Grund als hinreichend angesehen
werden kann, ohne dass eine vollständige Begründung vorliegt. Es wird
argumentiert, dass im Alltag regelmäßig Entscheidungen getroffen werden,
deren Begründung zwar nicht abschließend, aber ausreichend ist.
}

\begin{document}

\section{Einleitung}

Die Entscheidung für eine Handlung setzt nicht in jedem Fall eine
vollständige Kenntnis aller relevanten Umstände voraus.

In der Praxis genügt häufig ein Grund, der eine Entscheidung hinreichend
plausibel erscheinen lässt.

\section{Begriffliche Überlegungen}

Unter einem hinreichenden Grund soll im Folgenden ein Grund verstanden
werden, der für eine konkrete Entscheidung ausreicht, ohne sämtliche
möglichen Gegenargumente auszuräumen.

\section{Praktische Anwendung}

Besonders häufig findet sich dieses Prinzip bei alltäglichen Entscheidungen.
Die Frage, ob eine Entscheidung vollständig begründet werden kann, wird
dabei nicht immer gestellt.

\section{Schlussbemerkung}

Die Forderung nach vollständiger Begründung jeder alltäglichen Entscheidung
erscheint unter praktischen Gesichtspunkten nicht zwingend erforderlich.

Eine gewisse Restunsicherheit kann daher als Bestandteil einer
ordnungsgemäßen Entscheidungspraxis betrachtet werden.

\end{document}
